\documentclass[../main.tex]{subfiles}
\begin{document}
%==========================================TIÊU ĐỀ TẬP========================================
\begin{table}[h]
  \begin{adjustwidth}{-1cm}{}
    \begin{tabular}{>{\centering\arraybackslash}p{6cm}|>{\raggedright\arraybackslash}p{12.5cm}}
      \multirow{5}{6cm}
      % Cột bên trái: ÉPISODE và số tập
      {\centering
        {\\[-40pt]
          {\fontsize{20pt}{20pt}\selectfont{\outlinetext[1]{eptype}{white}
              {\flaregothic ÉPISODE}}\\[12pt]
            {\fontsize{72pt}{72pt}\selectfont{\outlinetext[1]{epnum}{white}
                {\burbank 156}}
            }}
          \\[8pt]
          \hrule
          \vspace{8pt}
          {\fontsize{20pt}{20pt}\selectfont{\outlinetext[1]{eptype}{white}
              {\flaregothic SPÉCIAL}}\\[12pt]
            {\fontsize{72pt}{72pt}\selectfont{\outlinetext[1]{epnum}{white}
                {\burbank XXX}}
            }}\\[6pt]
          {\fontsize{20pt}{20pt}\selectfont{\outlinetext[1]{epnum}{white}{\cafeteria (S-30)}}}\\[18pt]
          % {\fontsize{14pt}{14pt}\selectfont{\outlinetext[1]{epattr}{white}{\flaregothic [I]}}}\\
          {\fontsize{18pt}{18pt}\selectfont{\outlinetext[1]{epattr}{white}{\flaregothic COMPILATION}}}\\
          \color{black}
        }
      }
       &
      % Dòng thứ nhất: tên tập tiếng Nhật
      {\fotskip\fontsize{18pt}{18pt}\selectfont みんなが\ruby{注}{ちゅう}\ruby{目}{もく}したシーン\topten}                  \\[6pt]
      % Dòng thứ hai: tên tập tiếng Anh
       & {\cafeteria\fontsize{18pt}{18pt}\selectfont Highlight Scenes}              \\[4pt]
      % Dòng thứ ba: tên tập tiếng Việt
       & {\vnmsans\fontsize{18pt}{18pt}\selectfont Top 10 các cảnh đáng chú ý nhất}             \\[8pt]
      % Dòng thứ tư: Nhân vật xuất hiện
       & {\caviar{\fontsize{14pt}{14pt}\selectfont Nhân vật: \emp{sora1}}} \\
      % Dòng thứ năm (nếu có): Nhân vật chỉ xuất hiện trong Omake
       & {\abv Truyện phụ:} \emp{kuon2} \emp{kaigou2}
      \\[8pt]
      % Dòng cuối cùng: Thời gian diễn ra của tập (thường sẽ là ngày phát hành)
       & {\sen\fontsize{16pt}{16pt}\selectfont Mốc thời gian: 08/05/2022}  \\[18pt]
    \end{tabular}
  \end{adjustwidth}
\end{table}
%=========================================TRUYỆN CHÍNH========================================
\color{black}

\txtbx[2]{
  {\color{vol} \uline{Tóm tắt:}} Tập tổng hợp Top 10 cảnh đáng chú ý nhất của \textit{holo no Graffiti} trong 6 tháng, do Sora dẫn dắt dưới dạng hình tĩnh 2D nhằm ``cắt giảm kinh phí''. Trước đó vài ngày, Kuon kể lại cuộc bình chọn trên Twitter, ngạc nhiên vì cảnh AZKi ngậm micro lại đứng Top 1, đồng thời thắc mắc về khâu sản xuất ``hoạt hình'' kỳ quặc của tập này.
}

\begin{center}
  (Đây là danh sách top 10 do fan bình chọn cho những cảnh đáng chú ý nhất của \textit{holo no Graffiti} trong vòng sáu tháng qua.)
\end{center}

Khung hình mở ra với một sự cắt giảm kinh phí không thể lộ liễu hơn: Tokino Sora xuất hiện dưới dạng một tấm ``hình tĩnh'' 2D không hề cử động, nở một nụ cười tươi thắm rạng rỡ. Đằng sau cô là khung cảnh căn văn phòng làm việc quen thuộc nhưng cũng được vẽ bằng nét hoạt hình đơn giản đến mức buồn cười.

\img{s-30a}

Sora cất giọng tươi tắn, vang vọng như một MC chương trình truyền hình thực thụ: ``Đây là một tập điểm lại!''

Tấm hình tĩnh của cô nàng lập tức chuyển sang tư thế hai ngón tay trỏ chỉ thẳng xuống dưới, khóe miệng vẫn giữ nguyên nụ cười idol chuẩn mực: ``Tại sao chúng ta không nhìn lại \textit{hologra} trong vòng sáu tháng vừa qua nhỉ?''

Nói rồi, hình ảnh Sora tiếp tục đổi sang tư thế nháy mắt đầy tinh nghịch, đưa bàn tay trái ra phía trước thông báo đầy dõng dạc:

\begin{center}
  ``\textbf{`Top 10 cảnh của \textit{hologra} gây chú ý người xem' xin được phép bắt đầu!}''
\end{center}

Vẫn giữ nguyên dáng điệu idol tràn đầy năng lượng ấy, cô tiếp tục tung hứng: ``Chúng ta bắt đầu nhé! Bắt đầu với\dots''

\il

\bxh{10}{Khuôn mặt ác quỷ\dots}{153}{Mười Hai}    % S-29

[\dots] Vừa dứt câu thần chú, cả hai bản thể Polka, cùng với một chú Oruyanke nhỏ nhắn, hàng tá những bé Polka con tí hon chạy nhảy khắp nơi, và cả con slime xanh lè kia đều bị một lực hút khổng lồ kéo vào nhau, nhập lại thành một khối duy nhất. Sau một luồng ánh sáng chói lòa, hiện ra trước mắt mọi người là một khối dị hình không thể tả nổi: ``Polka Sushi'', một sự kết hợp kinh hoàng giữa hai con người và nguyên liệu làm sushi.

\gif{s-29q}{1}{155}

Game thảng thốt kêu lên, không thể chấp nhận cảnh tượng trước mắt: ``Hợp nhất lại đâu có phải như thế này đâu! Ai dạy cô hợp nhất kiểu điên rồ vậy!?''

Cả Kuon, Fubuki, Kaigou và Ryuuhou đều đứng hình, nhìn chằm chằm vào sản phẩm hợp thể kỳ quặc kia mà mồ hôi lạnh toát ra như mưa, đồng thanh thốt lên: ``Vậy cũng được sao? Cái cách hợp nhất này có hợp lý đâu?''

``Rồi bây giờ không biết gọi em ấy là Pol-Pol hay là Pol-Sushi nữa\dots'' cậu Tổng Kuon thở dài thượt, cảm thấy cuộc đời mình chưa bao giờ kỳ quái đến thế.

\begin{center}
  \uline{(Có một cái đồng hồ đếm ngược ở góc dưới màn hình. Đồng hồ chỉ về không khi cảnh ở dưới xảy ra.)}
\end{center}

Ngay khi hai người vừa hợp nhất lại thành khối Polka-Sushi dị hình, Okayu lập tức hai mắt sáng rực như đèn pha. Cô lao tới cục cơm bọc người đó, hét lên trong sự sung sướng tột cùng:

``Whoo! Cơm nóng kìa!''

Polka-cơm hoảng sợ kêu lên thất thanh rồi cuống cuồng bỏ chạy thục mạng.

Ryuuhou khoái chí reo hò, chạy vèo đuổi theo sau: ``Úi chà! Cơm nắm bọc người biết chạy kìa! Chờ em với Okayu-san ơi!''

Kaigou chạy ngay bên cạnh vội vã thò tay cắp nách kéo em gái lại, gắt gỏng: ``Em thôi xúi dại coi! Nó là con người chứ có phải đồ ăn đâu!''

Cả cậu Tổng, trợ lý Game và Fubuki cũng hoảng hốt đuổi theo phía sau. Cậu Tổng và Game đồng thanh thốt lên: ``Okayu-chan/Nekomata-san! Em/Cô ấy không phải là đồ ăn đâu!''

Fubuki dang hai tay cố gồng mình ngăn cản nàng Miêu nghiện cơm đang cố dùng sức cắn một miếng kia lại: ``Đừng! Bỏ cái đó ra đi, Oka-mew!''

Okayu run rẩy vì lên cơn thèm, đôi mắt đờ đẫn: ``Em thấy có gì đó trăng trắng\dots''

Mất hết lý trí, cô cắn mạnh một phát vào thứ ``trăng trắng'' đó. Nhưng trớ trêu thay, thứ cô cắn trúng lại chính là cánh tay của Fubuki!

``\textbf{Ái da\dots!!!}'' nàng Cáo trắng đau đớn gào lên một tiếng chói tai, đôi mắt lồi cả ra ngoài theo đúng nghĩa đen.

\begin{figure}[H]
    \centering
    \begin{subfigure}[t]{0.45\textwidth}
        \centering
        \includegraphics[width=8cm]{../../../../Image/Book12/s-29r1.png}
    \end{subfigure}
    \quad
    \begin{subfigure}[t]{0.45\textwidth}
        \centering
        \includegraphics[width=8cm]{../../../../Image/Book12/s-29r2.png}
    \end{subfigure}
\end{figure}
    
Kaigou nhăn mặt che mắt quay đi: ``Trời ơi nhìn thọt cả mắt ra ngoài kìa! Cắn tàn nhẫn vãi!'' Kuon đổ mồ hôi hột lao vào tìm cách gỡ nàng Mèo khỏi tay của nàng Cáo trắng: ``Okayu-chan, bình tĩnh lại đi!''

Game từ chiếc đồng hồ đeo tay cất giọng cảnh báo gấp giáp: ``\dots Không ổn rồi! Mọi người, giữ chắc vào!'' [\dots]

\bxh{9}{Sushi đủ loại hình dáng!}{154}{Mười Hai}    % 12-6

[\dots] Suisei đã xếp ra khay gỗ một miếng sushi đặc biệt. Đó là một Miko phiên bản chibi nhỏ xíu đang nằm sấp trên cục cơm trắng vắt tròn nóng hổi, thân mình được buộc chặt bằng một dải rong biển nori màu xanh đen.

``Của quý khách đây ạ.''

Miếng sushi Miko bất chợt ngẩng cái đầu tròn ủm lên nhìn Flare, đôi mắt rấn rấn như muốn thốt lên lời cảnh báo điều gì đó nhưng hoàn toàn bất lực, không thể cất thành tiếng. Nàng Bán yêu mỉm cười nhẹ nhàng cầm miếng sushi lên, chậm rãi đưa vào miệng rồi thong thả thưởng thức.

\img{12-6c}

\begin{center}
  \uline{(Có một cái đồng hồ đếm ngược ở góc dưới màn hình. Đồng hồ chỉ về không khi cảnh ở dưới xảy ra.)}
\end{center}

Thế nhưng, vừa nuốt chửng món ăn xong, Flare bỗng khựng lại, ngơ ngác cất tiếng: ``Hửm?''

*BỤP!*

Một luồng khói nhẹ tan ra, và quả nhiên điều kỳ lạ đã xảy ra. Cô lập tức biến thành một cục cơm nắm sushi nhỏ tẹo, với dáng nằm sấp hệt như món ăn mà cô vừa mới nuốt vào bụng. Giờ đây, cục sushi Flare đang nằm ngay ngắn trên chiếc khay gỗ, ngay bên cạnh Miko.

\img{12-6d}

Hai cục cơm nắm quay sang nhìn nhau. Đôi mắt nàng Vu nữ chớp chớp như muốn hét lên: ``Chị đã bảo rồi mà không chịu nghe!'', trong khi Flare mới ngơ ngác chợt nhận ra vấn đề. Cả hai đã hoàn toàn sập vào cái bẫy của tên chủ quán quỷ quyệt Hoshimachi.

Suisei thản nhiên bê chiếc khay gỗ cất xuống bên dưới quầy bar để làm mồi nhử cho vị khách tiếp theo, rồi quay trở lại lau dọn quầy bếp như chưa hề có chuyện gì xảy ra.

\bxh{8}{Mùi cá ngừ!}{134}{Mười Một}    % 11-1

[\dots] Bên trong tâm trí, ``Aki Tri Thức'' rùng mình ớn lạnh: ``Kinh quá\dots\ Cái trò gì thế này?''

\begin{center}
  \uline{(Có một cái đồng hồ đếm ngược ở góc dưới màn hình. Đồng hồ chỉ về không khi cảnh ở dưới xảy ra.)}
\end{center}

Bất chợt, một tiếng động lớn vang lên ngoài hành lang. Nene - trong vỏ bọc NeneRose đầy náo nhiệt - đột ngột lao xông qua cánh cửa văn phòng. Trên tay cô nàng Hải cẩu đang ôm khư khư một con cá cam khổng lồ vẫn còn tươi nguyên, hào hứng reo hò: ``Em câu được một con cá cam bự chảng nè!''

\img{11-1g}

``Aki Tri Thức'' hoảng hốt thốt lên: ``NeneRose!?''

Nhìn cô nàng Hải cẩu tự hào khoe chiến lợi phẩm, cả Kuon, Suzuno lẫn Game đồng thanh cất tiếng: ``Giờ không phải lúc khoe của đâu, Nene-chan / Nene-san/ Momosuzu-san!''

Bỏ ngoài tai những lời can ngăn, Nene vẫn tưng bừng ôm con cá nhảy tưng tưng quanh phòng, miệng liên tục hô vang: ``Cá cam! Cá cam! Cá cam!'' Trong khi đó, AkiRose ở bên ngoài lại tỏ ra vô cùng hưởng ứng, hai tay vỗ bộp bộp theo nhịp nhảy của đàn em.

Thế nhưng trong lúc hăng hái quá đà, Nene vô tình đưa con cá tươi rói lại quá gần mặt Aki. Dù Kuon và Suzuno vốn đã quen với những mùi hương nặng trong công việc hậu cần nên không gặp vấn đề gì, thì đối với khứu giác đang nhạy cảm vì say rượu của Aki, đó lại là một thảm họa thực sự.

\img{11-1h}

Trong tâm trí cô, câu hét ``\textbf{\textit{Ặc\dots\ mùi cá tanh quá\dots!}}'' vang lên đầy tuyệt vọng. [\dots]

\bxh{7}{Ai là thủ phạm?}{137}{Mười Một}    % 11-4

[\dots] Sora cúi gằm mặt xuống, im bặt không thốt ra một câu nào. Không gian phòng bỗng dưng lặng ngắt, tiếng thìa đũa cọ xát vào thành bát cũng tắt hẳn, chỉ còn lại tiếng thở nhẹ nhàng của mọi người đang lắng nghe. Ánh mắt tất cả đều đổ dồn vào bóng dáng nàng Thần tượng, sự hoang mang và lo lắng lan tỏa nhanh chóng như một làn sóng lạnh. Ai nấy đều đoán già đoán non, trong đầu dấy lên vô số suy nghĩ: có phải mì bị hỏng? Hay có vấn đề gì nghiêm trọng đang xảy ra? Bầu không khí căng thẳng dần dần bóp nghẹt mọi người, khiến mồ hôi lạnh bắt đầu xuất hiện trên trán một số thành viên.

\begin{center}
  \uline{(Có một cái đồng hồ đếm ngược ở góc dưới màn hình. Đồng hồ chỉ về không khi cảnh ở dưới xảy ra.)}
\end{center}

Phải một lúc lâu sau, cô mới ngẩng lên, phồng hai má nũng nịu dỗi: ``Có người trong phòng này đã tráo nước xúp của chị thành trà lúa mạch\dots'' Nói rồi, cô giận dỗi ngồi bệt xuống ghế.
Đến lúc này, mọi người trong phòng bắt đầu hỗn loạn đổ lỗi cho nhau.

Noel ngơ ngác quay sang nhìn Sora, đôi mắt tròn xoe không hiểu chuyện gì đang xảy ra: ``Ể?''
Fubuki cũng giật mình theo, lưỡi liềm trong tay còn đang dính dính mép bìa hộp: ``Thật à?''

Aki hoang mang lắc đầu, chiếc cọ lau bụi cứ thế buông thõng xuống sàn: ``Không thể nào\dots''
Fubuki vừa lau nhanh vết bụi còn sót trên tay, vừa hoảng hốt nhìn quanh phòng: ``Ai mà gan dám làm như thế chứ!?''

Game từ chiếc đồng hồ trên tay Kuon liền phán một câu xanh rờn, chẳng chút do dự: ``Cá chắc là Oozora-san luôn. Cô ấy là người bưng khay mì lên mà.''

Subaru nghe vậy giật nảy người, chiếc xe đẩy còn bị lắc mạnh đến phát ra tiếng kêu: ``Đừng đổ oan cho tôi chứ! Kaigou-san cũng đi bưng cùng tôi mà!''

Kaigou vội vã giơ cả hai tay lên trời, lắc đầu quầy quậy để phân trần: ``Chắc chắn không phải là em rồi! Em còn chả biết em mang trà lúa mạch về từ lúc nào nữa!''

Shion quay sang nhíu mày hỏi Miko, người đang đứng cạnh cô cạnh đó: ``Cậu nghĩ sao?'' trong khi đối phương chỉ biết gãi đầu cào tai, nhún vai chịu thua: ``Chịu. Không biết nữa.''

Aqua vội vàng xua tay, giọng nói lo lắng sợ bị dính oan: ``Không phải là tôi đâu nha!'' Luna cũng nhanh nhảu nối lời, chiếc kẹo mút trong miệng còn đang nhai ngấu nghiến: ``Cả em nữa-nora! Chắc là Fubuki-paisen làm rồi!''

Fubuki nghe vậy liền mếu xáo quay sang chỉ vào Luna, giọng oan ức không ai bằng: ``Ơ hay!? Chị còn chẳng dám đụng vào một sợi tóc của chị ấy nữa! Tin vào chị đi chứ!''

Miko đứng đó ngập ngừng, mắt nhìn quanh rồi lại dừng lại ở Kuon, định nói gì đó: ``Vậy thì\dots\ Kuon-sen--''

Sora lập tức ngắt lời Miko, tay đập nhẹ xuống bàn đầy chắc chắn: ``Chắc chắn không phải anh ấy đâu. Chị dám chắc anh ấy sẽ nói là: `Anh táng vào mồm bây giờ'.''

Kuon nghe vậy chỉ biết thở dài gãi đầu, giọng bất lực vang lên: ``Geez\dots\ Cảm ơn em vì đã nhắc.''

Towa cũng lên tiếng bênh vực Kuon, gật gù đồng tình với lời nói của Sora: ``Anh ấy còn thành thật hơn cả chị đấy, Vu nữ vụng về ạ.'' Miko nghe bị gọi là vụng về liền nổi đóa, chỉ tay vào nàng Ác ma mà gào lên: ``\textbf{Gọi ai là vụng về hả!?}''

Matsuri vội vàng bước ra can ngăn hai người, giơ tay xua đi cuộc cãi vã nhỏ sắp bùng nổ: ``Nào nào, hỏi xem ai đã thực sự làm thế đi.''

Fubuki không cần suy nghĩ gì thêm, liền chỉ tay thẳng vào Watame với quyết đoán không chút nghi ngờ: ``100\% là Watame làm luôn!'' Noel và Aki đứng bên cạnh liền gật đầu đồng tình tắp lưỡi, cùng nhau hô lên: ``Đồng ý\~{}''

Nàng Cừu thảng thốt ngơ ngác, chiếc hộp trong tay suýt nữa rơi xuống đất: ``Sao lại là em!?'' Game liền nhận xét thêm vào, giọng điệu còn khẳng định hơn mọi người: ``Nhìn mặt cô gian nhất trong cả nhóm đấy.''

Cô nàng nghe vậy liền khóc dở mếu dở, mắt long lanh nước sắp rơi: ``Nhưng mà em chẳng làm gì sai cả!!'' Towa đứng đó chỉ biết lắc đầu bất lực trước cảnh tượng hỗn loạn trong phòng: ``\dots Trời ạ.'' [\dots]

[\dots] Nhân tiện, nhìn lại khung cảnh cả hội ngồi dàn ngang bên bàn ăn lúc này\dots\ có ai thấy bức ảnh toàn cảnh ấy trông quen quen, giống hệt với kiệt tác ``The Last Supper'' (Bữa tối cuối cùng) của Leonardo da Vinci không nhỉ?

\img[18cm]{11-4i}

\bxh{6}{Watame-chan không phải là đồ ăn!}{136}{Mười Một}    % 11-3

[\dots] Chuyển sang câu hỏi thứ hai về ẩm thực, Botan sẽ là người thị phi. Towa dõng dạc: ``Câu hỏi số hai: ba món ăn được coi là cao lương mỹ vị hàng đầu trên thế giới là gan ngỗng, trứng cá tầm, và\dots?''

\begin{center}
  \uline{(Có một cái đồng hồ đếm ngược ở góc dưới màn hình. Đồng hồ chỉ về không khi cảnh ở dưới xảy ra.)}
\end{center}

Botan ranh mãnh nhếch môi: ``Watame chứ sao!''

\img{11-3g}

Nàng Sư tử trắng giật mở cái vung đậy thức ăn trên bàn ra, lộ ra một nàng Cừu đang co rúm sợ sệt. Botan nhanh như chớp rút sẵn bộ dao dĩa sáng bóng chuẩn bị ``thưởng thức'' tiền bối. Cả cậu Tổng lẫn trợ lý ảo đều đổ mồ hôi hột: ``Thịt cừu thì ở đâu mà chả có\dots''

Một lần nữa, câu trả lời này khiến ``Giáo sư'' Towa tức đến bốc khói đầu, nổi đóa: ``\textbf{Cậu ấy không phải là thức ăn!! Kuon-senpai, trả lời đi!}''

Kuon hắng giọng đáp: ``Nếu em đang nói về ẩm thực phương Tây thì món còn lại là nấm truffle, hay còn gọi là nấm cục.''

Towa đắc ý: ``Chính xác! Cái này chắc đã quá quen thuộc với mọi người rồi ha!''

Game lên tiếng từ đồng hồ: ``Tôi không nghĩ là ai cũng biết đâu\dots\ Tôi có dữ liệu về nước Pháp nên tôi mới biết đấy.''

Kuon gãi đầu: ``Tôi thì vẫn thích ăn phở hơn\dots''

Vị ``Giáo sư'' Ác ma hất cằm, cắt ngang lời của cậu: ``Câu tiếp theo!'' [\dots]

\bxh{5}{Hội ăn vỏ gối luộc\dots?}{123}{Mười}    % 10-3

[\dots] Bốn người nhìn tờ rơi. Rồi nhìn nhau. Sau đó lại nhìn về phía tờ rơi lần nữa.

Ryuuhou trố mắt chỉ vào tờ giấy: ``Khoan đã! Phí hội viên cho Chủ tịch lại là 100 nghìn tỷ yên!? Onii, dù anh có bán cả cái văn phòng này cộng với cả công ty thì chắc cũng không đủ tiền đóng phí cho cái hội kỳ quặc này đâu!''

``\dots Ừm\dots'' ba người còn lại cùng đồng thanh, ánh mắt hoang mang không thể nào che giấu được.

\begin{center}
  \uline{(Có một cái đồng hồ đếm ngược ở góc dưới màn hình. Đồng hồ chỉ về không khi cảnh ở dưới xảy ra.)}
\end{center}

Đúng lúc đó, tiếng kính vỡ vang lên lớn đanh gắt. Một cơn gió thoảng mạnh lùa vào phòng, ngay sau đó là một giọng nói đầy khí thế hô to: ``\textbf{Giờ bốn người đã biết bí mật rồi, thì không còn đường nào thoát khỏi Cái Chết Trắng đâu!}''

Flare và Roboco cùng đứng bật dậy trong trạng thái hoảng sợ: ``\textbf{Shirakami!?}''

Kuon cau chặt mày, cảm tưởng như vừa bị tạt một xô nước lạnh từ đầu xuống chân: ``Cái Chết Trắng\dots\ em đang nhắc đến \emph{ma túy} đấy à? Cái hội quái dị này rốt cuộc là cái gì vậy trời!?''

Vừa kịp nói xong câu đó, Fubuki (đã được thu nhỏ lại) lập tức tạo dáng như một nhân vật phản diện trong anime vừa mới lộ diện. Cô lao vội vào trong văn phòng (với tốc độ hình ảnh động giảm còn một nửa tới một phần ba so với thường lệ), tay tạo hình con cáo với dáng vẻ trông có vẻ ngầu nhưng lại vô cùng kệch cỡm.

\gif{10-3d}{1}{56}

Cô chỉ thẳng tay lên trần nhà và lớn tiếng tuyên bố: ``\textbf{Hội Thích Ăn Vỏ Gối của thành viên \hll\ như món luộc, hay còn gọi tắt là Hội Thích Ăn Vỏ Gối Luộc, là một tổ chức không ai có thể ngăn cản được! Hội ha-ha-ha!}''

Nàng Người máy nhăn mặt, chỉ vào tờ rơi trên bàn: ``Tên viết tắt của cái hội này gây hiểu lầm ghê thật, nghe kì cục lắm luôn.'' [\dots]

\bxh{4}{Gieo gió thì gặt bão thôi! (Tò mò cũng là một cái tội!)}{126}{Mười}    % 10-6

\begin{center}
  \uline{Xem lại {\flaregothic ÉPISODE 126}, quyển số Mười để đọc chú thich.}
\end{center}

\begin{center}
  [\dots] Lại thêm lũ mới lẻ loi,\\
  Là Thi Âm, Thỏ, với người Thiên phương.\\
  Chàng Viễn nhíu mày một đường,\\
  ``Lũ này lại tính tìm đường nghịch chi?''\\
  Phù thủy chạm khẽ thầm thì,\\
  Khuyển không động đậy, cả khi cựa mình.\\
  Họ cười, Bảo lắc đầu khinh,\\
  ``Khôn ngoan gì kiểu dọa thình lình ai?''\\
  Thi Âm liều lĩnh ra tay,\\
  Vừa cù đã bị đá bay ngang tường.
  
  \uline{(Có một cái đồng hồ đếm ngược ở góc dưới màn hình. Đồng hồ chỉ về không khi cảnh ở dưới xảy ra.)}

  Viễn cười: ``Anh đã lên giường-\\
  (À không), đã bảo đừng vương tay mà\dots [\dots]''

  \begin{figure}[H]
    \centering
    \begin{subfigure}[t]{0.45\textwidth}
      \centering
      \includegraphics[width=8cm]{../../../../Image/Book10/10-6f1.png}
    \end{subfigure}
    \quad
    \begin{subfigure}[t]{0.45\textwidth}
      \centering
      \includegraphics[width=8cm]{../../../../Image/Book10/10-6f2.png}
    \end{subfigure}
    \quad
    \begin{subfigure}[t]{0.45\textwidth}
      \centering
      \includegraphics[width=8cm]{../../../../Image/Book10/10-6f3.png}
    \end{subfigure}
  \end{figure}
\end{center}

{\color{bronze}\bxh{3}{Thác nước Polka!}{121}{Mười}}   % 10-1

[\dots] Polka chưa kịp đáp thì đôi mắt hình xiếc của cô bỗng sáng bừng lên, như vừa nghĩ ra trò gì đó khiến Kuon phải dè chừng: ``Em định giở trò gì đấy?''

Nàng Cáo sa mạc phớt lờ lời cậu, lập tức gọi tiền bối của mình, giọng tinh nghịch: ``\textbf{Mio-sen ơi!}''

Nàng Sói đen quay đầu lại, thấy Polka đang há miệng, nét mặt nũng nịu, như thể đang đòi được đút ăn. Nhưng thứ cô nhận lại chỉ là một cái nhìn trống rỗng từ Mio và một lời phản đối dứt khoát từ Kuon.

Cậu Tổng lắc đầu, giọng đầy mỉa mai: ``Mơ đi em ơi. Định há miệng chờ sung hả?'' Nàng Sói đen cũng không kém phần châm biếm: ``Cái gì thế, Merlion à?''

Ryuuhou bật cười trước phản ứng của hai người: ``Polka-san, chị mà muốn thành Merlion thì chị phải có nước đã. Để em đi lấy bình nước cho chị nhé?''

\begin{center}
  \uline{(Có một cái đồng hồ đếm ngược ở góc dưới màn hình. Đồng hồ chỉ về không khi cảnh ở dưới xảy ra.)}
\end{center}

Ngay lập tức, Polka bắt chước tượng Sư tử phun nước ở Singapore, miệng phun ra một dòng nước dài rõ rệt như một diễn viên hài sân khấu.

\begin{figure}[H]
  \centering
  \begin{subfigure}[t]{0.45\textwidth}
    \centering
    \includegraphics[width=8cm]{../../../../Image/Book10/10-1b1.png}
    \caption{Polka ra dấu hiệu muốn được nuông chiều\dots}
  \end{subfigure}
  \quad
  \begin{subfigure}[t]{0.45\textwidth}
    \centering
    \includegraphics[width=8cm]{../../../../Image/Book10/10-1b2.png}
    \caption{\dots nhưng cả Kuon và Mio không hiểu ý, kết quả là Polka làm trò cười cho thiên hạ.}
  \end{subfigure}
\end{figure}

Mio và Kuon đồng thanh phản ứng theo bản năng: ``Eo ơi!'' / ``Trời ơi, kinh quá!'' Ryuuhou ngồi trên ghế nghếch đầu lên xem, nhăn mặt cười xùy một tiếng: ``Trời ạ Polka-san, kĩ thuật phun nước này trặt nhịp hoàn toàn luôn ấy. Nhìn gớm chết đi được!''

Polka lau nước dãi trên miệng, mếu máo: ``Thôi nào, em chỉ muốn được anh chị đối xử như hoàng gia thôi\dots''

Kuon đảo mắt, tặc lưỡi một tiếng như đã đoán trúng ý định của cô: ``Đúng là đang ghen tị với Luna-chan rồi còn gì.'' [\dots]

{\color{silver}\bxh{2}{Xe máy chạy bằng cơm!}{132}{Mười}}   % 10-11

[\dots] Ayame nhún vai khoát tay, vẻ mặt điềm nhiên như chỉ vừa thấy con mèo nhà đi ngang qua phòng: ``Chuyện nhỏ thôi, để nàng ấy tự bình tĩnh lại đi. Cái quan trọng bây giờ là ta đã hoàn thành luyện tập thần thái, đã chuẩn bị gì cho bước cuối cùng chưa?'' cô quay sang hỏi hai đệ tử, mắt ánh lên sự háo hức không thể che giấu.

Nene và Watame cùng lúc gật gật đầu lia lịa, còn Ryuuhou ở bên cạnh thì nhíu mày, cảm thấy có gì đó không ổn sắp xảy ra. Ayame hất mái tóc ra sau, hô lớn làm cả phòng giật mình: ``Giờ đến bước cuối cùng rồi\dots\ \textbf{Đạp nẹt pô thôi mấy đứa!}''

``Ờ thôi, anh miễn nha,'' Kuon giơ tay từ chối ngay lập tức. ``Anh không gấu đến mức đó đâu.''

``Em cũng không tham gia đâu,'' Ryuuhou liền nối lời anh trai, lắc đầu quầy quậy. ``Em không muốn bị mời lên uống nước chè với các chú công an đâu ạ.''

``Ngươi thích thì ngươi có thể không tham gia được mà,'' nàng Quỷ sừng trả lời, nụ cười khẩy hiện rõ trên môi, chẳng hề bị lòng từ chối của hai người làm giảm nhiệt huyết.

Watame và Nene thì mắt sáng rực như trẻ con lần đầu thấy pháo hoa, lập tức nhảy cẫng lên: ``Nẹt pô thật hả chị?''

``Miễn là đừng làm gì quá điên rồ nha là được,'' Kuon nhắc khéo, giọng đã hơi lo lắng thực sự khi thấy ba cô gái đã sẵn sàng xông ra ngoài.

\begin{center}
  \uline{(Có một cái đồng hồ đếm ngược ở góc dưới màn hình. Đồng hồ chỉ về không khi cảnh ở dưới xảy ra.)}
\end{center}

\ct

Dưới ánh hoàng hôn tại bờ biển thành phố Điện tử, Kuon thong thả đạp chiếc xe đạp thuê, còn ba cô gái thì vừa đẩy khung xe máy tự chế, vừa chạy bộ trên đường mà hô vang những khẩu hiệu giang hồ.

\img{10-11g}

Nàng Quỷ sừng ngoái đầu lại, hét to về phía sau: ``\textbf{Mấy ngươi có thấy sướng không nào?}''

``C-Có ạ\dots'' nàng Cừu thở hổn hển đáp lời.

``T-Tuyệt vời quá\dots'' nàng Hải cẩu vừa thở dốc vừa cố nặn ra một nụ cười méo mó.

Cậu Tổng liếc qua, liền bật cười phì lên: ``Thích thế mà sao mặt mày mấy đứa lại tái nhợt ra thế?''

``Không\dots Không có gì đâu ạ\dots'' cả hai cùng đáp, cố gắng che giấu sự kiệt sức hiện rõ trên người.

Ayame cười lớn, vung tay hô to: ``\textbf{Được rồi! Tiến lên nào! Cùng nhau lao tới đường chân trời nào!}'' Bất ngờ có thêm sức lực từ đâu, Watame và Nene lập tức đồng thanh hét lên như vừa nhận được lệnh từ người thủ lĩnh của mình. [\dots]

{\color{gold}\bxh{1}{AZKi xuất hiện lần đầu!}{145}{Mười Một}}   % 11-11

[\dots] Watame lúc này gần như không thở được, hốt hoảng gào lên trong hơi thở thều thào. Cô cầu xin được tha thứ: ``Cho em\dots\ cho em một cơ hội nữa đi mà!! Lần này em sẽ làm chuẩn cho mọi người xem!!''

Kuon nhanh chóng tiến lại gần để can ngăn, cố gắng gỡ tay Subaru ra khỏi cổ Watame. Cậu nói với giọng cầu xin: ``Còn bài vở nào em cứ tung hết ra đi Watame-chan. Mà Suba-chan này, bỏ em ấy ra đi, kẹp thế kia là nghẹt thở đi chầu ông bà bây giờ đấy.''

Cuối cùng, phải nhờ sự phối hợp của cả Kuon, Kaigou và AZKi cùng nhau xắn tay áo kéo nàng Vịt ra, Watame mới may mắn thoát khỏi lưỡi hái tử thần trong tích tắc.

\begin{center}
  \uline{(Có một cái đồng hồ đếm ngược ở góc dưới màn hình. Đồng hồ chỉ về không khi cảnh ở dưới xảy ra.)}
\end{center}

\ct

\begin{center}
  \uline{Đoạn dưới đây diễn ra trong thước phim quay thử.}
\end{center}

``[\dots] Thôi chết rồi! Muộn giờ học mất rồi!!'' AZKi hốt hoảng vừa chạy vội vã dọc hành lang vừa thốt lên, diễn tròn vai nữ sinh muộn học.

Thế nhưng, thay vì lát bánh mì gối ban đầu, thứ đang bị kẹp chặt giữa hai hàm răng của cô nàng lại là\dots\ một chiếc micro cầm tay bằng kim loại sáng loáng, bóng lộn đến lạ.

\img{11-11j}

Cả Kuon đang đứng xem và Game (qua loa máy tính) lập tức giật khóe mắt, cảm giác déjà vu quen thuộc ùa về như đã từng thấy cảnh này ở đâu đó.

Ryuuhou gãi đầu bối rối, không hiểu nổi tình hình: ``Sao en thấy cái kịch bản này nó quen quen thế nhỉ\dots''

Giọng trợ lý ảo lập tức chệm vào từ loa, giọng đầy lo lắng: ``Đừng nói với tôi là chúng ta định bắt chước cái môtíp cliché anime kiểu nữ sinh ngậm bánh mì chạy vội, va vào nam chính ở góc đường rồi quay ra trở thành người yêu đấy nhé!?''

Vừa dứt câu, cả Kaigou lẫn Watame đồng loạt quay sang nhìn Kuon với gương mặt vênh lên đầy ẩn ý, như thể vừa bị nói trúng tim đen.

Cậu Tổng tức thì nhăn nhó mặt mày, hét lên không hiểu nổi: ``\textbf{Mấy người nhìn tôi làm cái quái gì!?}'' Cậu hắng giọng một tiếng rồi chỉ vào màn hình máy quay: ``Với cả\dots\ tại sao cái lát bánh mì lúc nãy lại biến thành cái micro rồi!?''

Game thản nhiên giải thích như thể đó là điều hiển nhiên nhất: ``Tôi vừa chỉnh sửa hiệu ứng thực tế ảo đôi chút trên máy. AZKi-san mở lời nhờ tôi làm thế đấy.''

Kaigou liền gật gù tán thành, đồng tình với ý tưởng độc đáo đó: ``Đổi thành cái micro cho nó hợp với hình tượng ca sĩ mê hát của cậu ấy chứ sao!''

Trong thước phim, AZKi vẫn tiếp tục chạy hối hả, vừa thở hổn hển vừa cất giọng giới thiệu bản thân với khán giả: ``Mình là AZKi-chi! Micro tụ điện (condenser mics) chính là nguồn sống và chân lý đời mình!!'' Câu giới thiệu kỳ dị trên khiến Kuon đứng ngoài cảm thấy rùng mình ớn lạnh, còn Ryuuhou thì toát mồ hôi hột vì không biết chuyện sẽ đi về đâu.

Nàng Trưởng nhóm đổ mồ hôi hột với cách xây dựng tình huống: ``Khán giả xem xong sẽ tưởng chị có sở thích quái dị là đi nhai thiết bị âm thanh đấy, AZKi-san\dots''

Kaigou cười gượng gạo, cố gắng bao biện cho tình thế khó xử: ``Yêu thích ca hát đến mức ăn cả micro luôn kìa\dots''

Bất ngờ, giọng của AZKi trong phim đổi sang hiệu ứng vọng âm (reverb) đầy huyền bí, như một lời tiên tri: ``Nhưng loại micro mà mình đang ngậm trong miệng đây lại là Shure SM58! Thuộc dòng micro điện động (dynamic mic) nhé!!''

Game cất giọng cạn lời, hoàn toàn bó tay trước tình hình: ``Cô nói mấy cái thuật ngữ chuyên ngành đó thì người xem bình thường chẳng ai hiểu đâu!''

Nghe vậy, Ryuuhou liền gật gù phân tích cực kỳ nghiêm túc như một chuyên gia âm thanh thực thụ: ``Nhưng mà Azu-san nói đúng đấy! Dòng Shure SM58 là `huyền thoại' sống của các buổi diễn live nhờ độ bền khủng và chống hú tốt.'' Cô ấy dừng lại một chút rồi bổ sung: ``Có điều ngậm cái đầu củ màng bọc kim loại đó trong miệng thì ê răng chết đi được\dots''

Đúng lúc này, trên màn hình bất ngờ xuất hiện hình vẽ hoạt họa của một vị ``Giáo sư âm nhạc tự xưng'', mặc áo choàng đen và đeo kính cận lớn. Vị này giọng hấn hắng phát biểu với gương mặt đầy tự mãn: ``Chắc hẳn các bạn đang tự hỏi micro tụ điện khác gì với micro điện động đúng không? Vậy để tôi giải đáp thắc mắc chuyên sâu cho các bạn nhé--''

*RẮC!*

Chưa kịp để vị giáo sư hoạt họa ấy thao thao bất tuyệt, Subaru từ đâu lao tới như một cơn lốc, tay xé toạc màn hình hoạt họa vừa xuất hiện. Cô tức giận gào lên bằng 100\% công suất thanh quản, tiếng vọng đi khắp hành lang: ``\textbf{Méo ai quan tâm mấy cái này đâu bà nộiii!!}'' [\dots]

\il

Sau một hồi dài điểm lại toàn bộ những khoảnh khắc ``khó đỡ'' ấy, khung cảnh quay trở lại với gian văn phòng hoạt hình ban đầu. Sora buông một tiếng thở dài thượt, giọng điệu có phần mệt nhoài sau khi phải gánh toàn bộ phần thuyết minh cho mớ hỗn độn vừa rồi: ``Haiya\~{} Các cảnh ấy thật điên rồ phải không nào?''

Hình ảnh của cô nàng ``xoay'' nhẹ một góc, quay trở lại dáng chỉ hai tay trỏ xuống dưới: ``Nếu có bất cứ cảnh nào làm cho các bạn thích thú, các bạn có thể xem lại các tập đó ở đầu mỗi phần nhé!''

Sora chắp hai tay lại, mỉm cười gửi lời nhắn nhủ chân thành đến khán giả: ``Tuần sau, chúng ta sẽ tiếp tục với \textit{hologra} như bình thường nhé! Mong các bạn sẽ ủng hộ!''

Để kết thúc thước phim, hình ảnh Sora chuyển sang tư thế tự đưa hai nắm đấm nhỏ nhắn lên áp sát đầu\dots\ Có lẽ cô nghĩ làm như thế trông sẽ vô cùng đáng yêu chăng?

\img{s-30b}

Ừ thì, cứ cho là như vậy đi.

%==========================================TRUYỆN PHỤ=========================================
\omk

\pov{Hikawa Kuon}

\begin{center}
  \uline{Hồi tưởng\\Vài ngày trước.}
\end{center}

Vì tập \textit{hologra} tuần này sẽ là một tập tổng hợp Top 10, nên A-san và tôi đã cùng nhau lập ra một cuộc thăm dò ý kiến trên Twitter với câu hỏi: ``Cảnh nào khiến cho các bạn tâm đắc nhất? Và ở tập nào?''

Chỉ trong vòng ít giờ sau khi đăng tải, chúng tôi đã nhận được hàng tá câu trả lời gửi về. Sự nhiệt tình này không làm tôi quá ngạc nhiên, bởi bản thân tôi khi xem lại những thước phim ấy đôi lúc cũng cảm thấy khá thú vị.

Ngoại trừ một chi tiết duy nhất. Cảnh quay xuất sắc giật giải quán quân Top 1 lại chính là phân đoạn khi AZKi lần đầu tiên bước chân đến văn phòng của chúng tôi.

Tôi ngồi tựa lưng vào ghế làm việc, tay lật lật bản danh sách kết quả vừa in ra, quay sang nhấp một ngụm cà phê rồi cất giọng nói với Kaigou đang đứng pha trà gần đó: ``Cái cảnh mà chúng ta quay phim em ấy chạy với cái míc-rô ngậm trên miệng\dots'' tôi chẹp miệng cảm thán.

Em ấy dừng tay pha trà, khẽ nghiêng đầu hồi tưởng rồi gật gù: ``À\~{} Cái cảnh đó. Nó cũng rất là kỳ quặc luôn.''

Tôi cười trừ một tiếng đầy ngao nán: ``Ờ. Chả hiểu vì sao người ta lại chọn cái đó nữa. Lúc anh đưa bảng kết quả cho AZKi-chan xem, mặt em ấy đỏ bừng lên như trái cà chua chín, cứ ôm mặt than `Sao mọi người lại nhớ tới cái cảnh xấu hổ đó chứ!' ''

Kaigou đặt đĩa tách trà xuống bàn, khẽ mỉm cười nhếch môi: ``Thì chính cái sự đối lập giữa hình tượng diva thanh lịch của AZKi và cái cảnh ngậm micro chạy vắt chân lên cổ đó mới tạo nên sức hút mà. A-chan lúc nhìn thấy kết quả chắc cũng cười ngất đúng không?''

``Còn phải hỏi,'' tôi thở dài, đưa tay lên day day thái dương. ``Chị ấy còn định dùng đoạn đó làm hình động đăng lên tài khoản chính thức cơ, may mà anh cản lại kịp. Nhưng thôi, ít nhất trong bảng xếp hạng cũng không có cảnh nào quá vô lý hoặc quá hở hang để bị đưa vào Top 10. Thế cũng coi như an ủi rồi.''

Tôi tựa người ra sau ghế, liếc nhìn vào màn hình máy tính đang chiếu thử bản dựng thô của tập tuần này, rốt cuộc vẫn không kiềm được mà chêm thêm một câu: ``Cơ mà, anh vẫn không thể nào nuốt nổi khâu sản xuất của tập này. Kiểu, thật sự nghĩ rằng bên phía biên tập cũng phải đầu tư lấy một phông nền nào đó tử tế hơn chứ? Hoặc chí ít thì họ hoàn toàn có thể quay trực tiếp ngay tại văn phòng này cơ mà? Hà cớ gì mà phải dày công tạo ra một cái nền văn phòng vẽ tay hoạt hình đơn điệu, rồi lại để Sora-chan làm một tấm hình tĩnh 2D đứng thuyết minh cơ chứ?''

Kaigou khoanh tay tựa vào cạnh bàn, đôi mắt xanh lấp lánh sự trêu chọc nhìn tôi: ``Này Aniki, chẳng phải tuần trước chính anh là người đã ký duyệt cắt giảm ngân sách truyền thông cho tập tổng hợp sao?''

Tôi khựng lại một nhịp, suýt chút nữa thì sặc ngụm cà phê trong miệng. Tôi vội vàng lên tiếng bào chữa: ``Này nhé! Anh duyệt cắt giảm ngân sách hậu kỳ để dồn tiền cho dự án sân khấu tháng sau, chứ đâu có bảo bọn nó xài cái tấm hình PNG không cử động của Sora-chan rồi vẽ cái phông nền trông như vẽ bằng phần mềm Paint thế kia!''

Cô em gái nhấp một ngụm trà, mỉm cười nhẹ nhàng: ``Nhưng công nhận Sora-chan giỏi thật đấy. Chỉ với một tấm hình tĩnh không nhúc nhích mà em ấy vẫn tung hứng, diễn xuất tràn đầy năng lượng như một idol chuyên nghiệp trên sóng truyền hình vậy. Khán giả bên dưới phần bình luận thậm chí còn khen đây là cái tập `nhảm nhí một cách nghệ thuật' nhất từ trước đến nay nữa kìa.''

``Đó chính là cái anh lo lắng đấy,' tôi khoanh tay lắc đầu ngán ngẩm trước cái tư duy logic không giống ai của cái công ty này. ``Cứ cái đà được đà xông lên này thì lần sau bên biên tập sẽ cắt luôn cả kinh phí lồng tiếng, rồi cho nhân vật giơ mấy cái bảng chữ tay lên mất thôi.''

Kaigou bật cười khẽ, âm thanh trong trẻo xua tan đi sự căng thẳng trong gian phòng: ``Nếu chuyện đó xảy ra thật thì anh lại phải đứng ra thu dọn hậu quả thôi. Dù sao thì\dots\ `vì \hll\ là thế mà'.''

Tôi thở dài một hơi, nhìn cô nàng rồi lại nhìn sang bản danh sách phản hồi tràn ngập tiếng cười của người hâm mộ trên màn hình.

Mà thôi, thấy mọi người vui vẻ như vậy, tôi cũng chẳng buồn chấp nhặt làm gì. Không hề.
\end{document}